\documentclass[11pt,a4paper]{article}

\usepackage{../shared/cathedral-preamble}
\usepackage{setspace}
\onehalfspacing
\definecolor{correct}{RGB}{0, 128, 0}
\definecolor{wrong}{RGB}{200, 0, 0}

\title{%
  Reporting on the Cathedral:\\
  A Guide for Science Journalists%
}
\author{
  Jason Robert Gochanour
}
\date{April 20, 2026}

\begin{document}

\maketitle

\begin{abstract}
This document provides accurate framing for journalists covering
the Cathedral project: a machine-verified reduction of the
Riemann Hypothesis in Lean~4. It contains: (1)~what the project
is, in one paragraph; (2)~what it is \emph{not}, which is more
important; (3)~a fact sheet with verified numbers; (4)~common
misconceptions and their corrections; (5)~suggested expert sources
for independent evaluation; and (6)~quotes approved for use.
The primary goal of this document is to prevent misreporting.
The Riemann Hypothesis has \textbf{not been proved}. Cryptography
has \textbf{not been broken}. What has been built is significant
in its own right, and accurate reporting will serve the public
better than sensationalism.
\end{abstract}

% ================================================================
\section{The One-Paragraph Summary}
% ================================================================

A computer scientist in New Mexico, working with AI assistance over
30 days, built a 161-file computer program in the Lean~4 theorem
prover that establishes, with machine-verified certainty, that the
167-year-old Riemann Hypothesis is logically equivalent to the
convergence of a concrete mathematical approximation problem. The
proof is complete in both directions and relies on exactly two
mathematical assumptions (``axioms''), both of which are well-
understood results in classical mathematics that have not yet been
translated into computer-checkable form. The project does not prove
the Riemann Hypothesis. It provides the most precise formal map ever
made of what a proof would require.

% ================================================================
\section{What It Is NOT}
\label{sec:not}
% ================================================================

This section is more important than the previous one.

\begin{enumerate}
  \item \textcolor{wrong}{\textbf{NOT a proof of the Riemann
    Hypothesis.}} The project reduces RH to two axioms. It does
    not prove those axioms. The distinction is critical: a
    reduction is a map of the problem, not a solution to it.

  \item \textcolor{wrong}{\textbf{NOT a threat to cryptography.}}
    A companion paper (\emph{cathedral-security.tex}) establishes
    rigorously that the mathematical tools used in the Cathedral
    provide no computational advantage for breaking encryption.
    Even a full proof of RH would not break RSA, AES, or
    post-quantum cryptography. Headlines like ``Riemann Hypothesis
    breakthrough threatens internet security'' would be factually
    wrong.

  \item \textcolor{wrong}{\textbf{NOT AI solving mathematics
    autonomously.}} The AI generated code under human direction.
    The human provided architectural vision, quality control, and
    mathematical taste. The AI did not ``solve'' anything---it
    implemented strategies designed by the human and verified by
    the Lean compiler.

  \item \textcolor{wrong}{\textbf{NOT a claim to the Millennium
    Prize.}} The Clay Mathematics Institute requires a proof of
    RH itself, not an equivalence. This project would not qualify.

  \item \textcolor{wrong}{\textbf{NOT a professional mathematics
    paper.}} The author is a computer scientist, not a
    professional mathematician. The project is a formalization
    (software engineering applied to mathematics), not a
    traditional research publication.
\end{enumerate}

% ================================================================
\section{What It IS}
% ================================================================

\begin{enumerate}
  \item \textcolor{correct}{\textbf{A machine-verified formal
    proof.}} Every logical step has been checked by the Lean~4
    compiler. There are zero errors, zero gaps, and zero
    unchecked assumptions beyond the two stated crown axioms.

  \item \textcolor{correct}{\textbf{A complete biconditional.}}
    The equivalence goes in both directions: RH implies convergence
    (forward), and convergence implies RH (converse). Both are
    verified.

  \item \textcolor{correct}{\textbf{A case study in human-AI
    collaboration.}} One person with AI assistance produced in
    30 days what would traditionally require a research team over
    months. This has implications for how mathematical research
    is conducted.

  \item \textcolor{correct}{\textbf{A transparent artifact.}}
    All 161 active files, 128 archived files, 55 axioms, and
    theorems are available for inspection. Everything is
    documented.
\end{enumerate}

% ================================================================
\section{Fact Sheet}
% ================================================================

\begin{center}
\begin{tabular}{@{}ll@{}}
\toprule
\textbf{Item} & \textbf{Verified Value} \\
\midrule
Active Lean files & 161 \\
Archived Lean files & 128 \\
Total axioms (active) & 55 \\
Crown axioms (on critical path) & 2 \\
Machine-verified theorems & $\sim$1,335 \\
Compilation errors & 0 \\
Development time & 30 days \\
Team size & 1 human + AI assistants \\
Prover & Lean 4 \\
Mathematical library & Mathlib \\
Companion papers & 21 \\
Millennium Prize eligibility & No \\
Cryptographic impact & None \\
\bottomrule
\end{tabular}
\end{center}

% ================================================================
\section{Common Misconceptions --- Corrected}
% ================================================================

\begin{center}
\small
\begin{tabular}{@{}p{6cm}p{6.5cm}@{}}
\toprule
\textcolor{wrong}{\textbf{Wrong}} &
\textcolor{correct}{\textbf{Right}} \\
\midrule
``The Riemann Hypothesis has been proved.'' &
The RH has been \emph{reduced} to two well-understood
axioms. It has not been proved. \\[6pt]

``This breaks encryption.'' &
A companion security paper proves it does not. Even
a full proof of RH would not break modern cryptography. \\[6pt]

``AI solved a famous math problem.'' &
A human used AI as a coding assistant to build a formal
verification system. The Lean compiler verified it. \\[6pt]

``This is just a computer program, not real math.'' &
Machine-verified proofs are \emph{more} rigorous than
traditional proofs: every step is checked, zero gaps. \\[6pt]

``The project uses 40 unproved assumptions.'' &
55 axioms total, but only 2 are on the crown theorem's
critical path. The other 53 support alternative proof
routes. \\[6pt]

``This was done in 30 days, so it can't be serious.'' &
The tools (Lean 4, Mathlib, AI) have transformed what
is possible. The result is 1,335 verified theorems, not
a hasty sketch.  \\
\bottomrule
\end{tabular}
\end{center}

% ================================================================
\section{The Riemann Hypothesis --- Explained Simply}
% ================================================================

For readers unfamiliar with the mathematics:

\textbf{Prime numbers} (2, 3, 5, 7, 11, 13, \ldots) are the
building blocks of all whole numbers. Every number can be written
as a product of primes. We know roughly how many primes there are
below any given number, but we don't know their exact pattern.

In 1859, Bernhard Riemann found a mathematical function (the
``zeta function'') that encodes where the primes are. This function
has special values called ``zeros'' that control the pattern of
primes. Riemann conjectured that all these zeros line up perfectly
on a single line in the complex plane.

If true, primes are as orderly as we hope. If false, there is
deep irregularity in the fabric of arithmetic.

The Cathedral proves that this conjecture is \emph{exactly
equivalent} to whether certain wiggly curves (sawtooth waves) can
be combined to make a flat line. Both directions of this equivalence
are machine-verified.

% ================================================================
\section{Suggested Expert Sources}
% ================================================================

For independent evaluation, journalists may wish to consult:

\begin{itemize}
  \item \textbf{Lean/Mathlib community}: The Lean Zulip chat
    (\url{https://leanprover.zulipchat.com/}) includes experts
    who can evaluate the formal proof.

  \item \textbf{Analytic number theory}: Any university faculty
    member specializing in analytic number theory can evaluate
    whether the two axioms are standard results.

  \item \textbf{Formal verification}: Experts at Carnegie Mellon,
    Imperial College, or Microsoft Research can evaluate the
    verification methodology.

  \item \textbf{AI and mathematics}: Google DeepMind (AlphaProof
    team) or equivalent groups can evaluate the human-AI
    collaboration claims.
\end{itemize}

% ================================================================
\section{Approved Quotes}
% ================================================================

The following quotes may be attributed to the author:

\begin{quote}
``We did not prove the Riemann Hypothesis. We built a machine-
verified map of exactly what a proof would require. The map has
two blank spots. Those blank spots are the hardest parts.''
\end{quote}

\begin{quote}
``The Lean compiler doesn't care who you are. It only cares
whether your proof is correct. That's what makes machine
verification powerful: it's the most objective form of
mathematical truth we have.''
\end{quote}

\begin{quote}
``This project does not break cryptography, does not solve a
Millennium Prize Problem, and does not mean AI is replacing
mathematicians. What it does mean is that one person with the
right tools can contribute to frontier mathematics. That's
new.''
\end{quote}

% ================================================================
\section{What NOT to Write}
% ================================================================

Please avoid:

\begin{itemize}
  \item Headlines containing ``proved,'' ``solved,'' or ``cracked''
    in reference to the Riemann Hypothesis.
  \item Any implication that encryption or internet security is
    affected.
  \item Framing the AI as an independent agent that ``discovered''
    mathematical results.
  \item Comparisons to the Millennium Prize that imply eligibility.
  \item Use of the word ``breakthrough'' without the qualification
    ``in formalization methodology.''
\end{itemize}

Accurate alternatives:

\begin{itemize}
  \item ``Machine-verified reduction of the Riemann Hypothesis''
  \item ``Formal proof maps famous conjecture to two axioms''
  \item ``Human-AI collaboration produces verified math in 30 days''
  \item ``Computer checks every step of 167-year-old problem's
    logical structure''
\end{itemize}

% ================================================================
\section{Contact}
% ================================================================

For questions, clarifications, or interview requests, contact
the author directly. All companion papers are available upon
request.

\end{document}
